% Every dimension and part for the chess robot, in one place.
% Build with:  tectonic dimensions.tex
% The numbers come from robotics/cad/src/lib/dims.py. If you change one there,
% change it here too.

\documentclass[10pt]{article}

\usepackage[margin=0.85in]{geometry}
\usepackage[T1]{fontenc}
\usepackage{lmodern}
\usepackage{microtype}
\usepackage{booktabs}
\usepackage{tabularx}
\usepackage{array}
\usepackage{enumitem}
\usepackage{xcolor}
\usepackage[hidelinks]{hyperref}
\usepackage{titlesec}
\usepackage{needspace}

\definecolor{accent}{HTML}{1F4E79}
\definecolor{muted}{HTML}{555555}
% needspace keeps a heading on the same page as the table under it
\titleformat{\section}{\needspace{14\baselineskip}\large\bfseries\color{accent}}{\thesection}{0.6em}{}
\titleformat{\subsection}{\needspace{10\baselineskip}\normalsize\bfseries}{\thesubsection}{0.6em}{}
\titlespacing*{\section}{0pt}{1.4em}{0.5em}
\titlespacing*{\subsection}{0pt}{1.0em}{0.3em}

\setlength{\parindent}{0pt}
\setlength{\parskip}{0.45em}
\setlist{nosep, leftmargin=1.4em}
\renewcommand{\arraystretch}{1.15}

\newcolumntype{L}{>{\raggedright\arraybackslash}X}
\newcommand{\mm}{\,mm}
\newcommand{\todo}[1]{\textcolor{red!70!black}{\textbf{To decide:} #1}}
\newcommand{\file}[1]{\texttt{#1}}

\title{\color{accent}\textbf{Chess robot: dimensions and parts}}
\author{Robotics subsystem, Tufts senior design}
\date{November 27, 2026}

\begin{document}
\maketitle
\vspace{-1.5em}

\begin{center}
\begin{minipage}{0.9\textwidth}
\small\color{muted}
All sizes are in millimeters. The CAD reads its numbers from \file{robotics/cad/src/lib/dims.py},
so that file wins if it ever disagrees with this one. Anything in \textcolor{red!70!black}{\textbf{red}}
hasn't been designed or bought yet.
\end{minipage}
\end{center}

\setcounter{tocdepth}{1}
\tableofcontents

% ---------------------------------------------------------------------------
\section{How to read the numbers}

There are two reference planes.
\begin{itemize}
  \item \textbf{Pieces:} $z = 0$ is the \emph{contact plane}, the bottom of the brass washer and the
        bolt head. The plastic starts 0.3\mm{} higher, so only metal touches the board.
  \item \textbf{Claw:} $z = 0$ is the \emph{seat plane}, where the top of a piece's cap cone sits
        while the piece hangs in the claw. Negative $z$ is below it, toward the board.
  \item \textbf{Gantry:} \emph{board coordinates}. $x = y = 0$ is the center of the board and $z = 0$
        is the top of the pads. $x$ runs along the runways.
\end{itemize}
``r'' means radius and ``\O'' means diameter.

% ---------------------------------------------------------------------------
\section{Board}

\begin{tabularx}{\textwidth}{l l L}
\toprule
Item & Size & Notes \\
\midrule
Square & 50 & FIDE allows 50--65. The smallest size keeps the gantry small \\
Playing area & 400 $\times$ 400 & 8 squares of 50 \\
Border & 25 & On every side \\
Whole board & 450 $\times$ 450 & Too big for one print. Printed tiles or a laser-cut plate, to be decided after the 2$\times$2 test \\
\bottomrule
\end{tabularx}

\subsection{Contact pads on each square}

Each square gets two copper tape pads. With these sizes a piece can sit up to 3\mm{} off
center in any direction and still read correctly without shorting. (It was 4\mm{} with the
bigger M12 washer. The M10 washer's smaller hole costs the millimeter.)
\file{checks/fit\_check.py} checks this.

\begin{tabularx}{\textwidth}{l l L}
\toprule
Pad & Size & What it touches \\
\midrule
Center dot & r 1.5 (\O{} 3.0) & The M3 bolt head \\
Outer ring & r 6.5 to 13.0 (\O{} 13 to 26) & The M10 brass washer \\
Placement tolerance & $\pm$3.0 & How far off center a piece can be \\
\bottomrule
\end{tabularx}

Worst-case margins at 3\mm{} off center: center dot to washer 0.75, ring to bolt head 0.5,
ring still on the washer by 0.5, center dot still on the bolt head by 1.25. The ring needs copper
tape at least 26\mm{} wide, or two strips side by side.

\subsection{2$\times$2 test tile (\file{board/tile\_2x2})}

\begin{tabularx}{\textwidth}{l l L}
\toprule
Feature & Size & Notes \\
\midrule
Footprint & 100 $\times$ 100 & Four 50\mm{} squares \\
Top plate & 4 thick & The pad surface is the top of this plate \\
Rails underneath & 4 wide, 6 tall & Room for wiring under the plate \\
Total height & 10 & \\
Engraved guide lines & 0.6 wide, 0.4 deep & Square borders and the three pad circles, for cutting the tape \\
Wire holes & \O{} 1.6 & Two per square: one at the center, one 11.75 from center (just outside the steel washer) \\
Steel washer pocket & \O{} 20.3, from underneath & Stops 1.0 below the pads. Glue an M10 steel washer into each one \\
\bottomrule
\end{tabularx}

Print the tile upside down so the pad side is on the bed and comes out flat. The washer pockets
then face up and need no supports.

% ---------------------------------------------------------------------------
\section{Pieces}

Heights are the FIDE tournament set scaled by 0.88 (our 50\mm{} squares vs.\ their 57).
The king's base is 41\% of its height, and FIDE asks for 40--50\%.

\begin{tabularx}{\textwidth}{l c c c L}
\toprule
Piece & Height & Base \O & Plastic at 100\% infill & Notes \\
\midrule
Pawn   & 45 & 28 & $\approx$12\,g & \\
Rook   & 50 & 30 & $\approx$16\,g & \\
Knight & 55 & 30 & $\approx$17\,g & 12\mm{} thick head profile \\
Bishop & 62 & 30 & $\approx$16\,g & \\
Queen  & 75 & 33 & $\approx$21\,g & \\
King   & 83 & 34 & $\approx$25\,g & Cross is on the front, since the top is the grab point \\
\bottomrule
\end{tabularx}

The washer, bolt, nut and two magnets add about 7\,g to every piece, so the heaviest piece (the king) is about 32\,g.

\subsection{Body sections shared by every piece}

\begin{tabularx}{\textwidth}{l l L}
\toprule
Section & Size & Notes \\
\midrule
Gap under the plastic & 0.3 & The metal sticks out this far below the plastic \\
Foot & 7 tall, base \O{} & 0.6 chamfer at the bottom (no first-layer squish), 1.0 at the top. Tall enough to wrap around the magnets \\
Stem & r 9, starts at $z = 16$ & The foot tapers into this \\
\bottomrule
\end{tabularx}

\subsection{Mushroom head (identical on every piece)}

This is the only part of a piece the claw ever touches. With $H$ = piece height:

\begin{tabularx}{\textwidth}{l l l L}
\toprule
Part & Size & From $z$ to $z$ & Notes \\
\midrule
Cap flat top & r 12 (\O{} 24), 2 tall & $H-2$ to $H$ & \\
Cap cone & r 12 down to r 5, 7 tall & $H-9$ to $H-2$ & 45\textdegree, prints without support. Seats in the jaws \\
Neck & r 5, 5 tall & $H-14$ to $H-9$ & The jaws close here. Nothing on the body may go past r 5 in this band \\
Whole head & 14 tall & $H-14$ to $H$ & \\
\bottomrule
\end{tabularx}

% ---------------------------------------------------------------------------
\section{Piece contacts and hardware}

\subsection{Bought parts, per piece}

\begin{tabularx}{\textwidth}{c L l L}
\toprule
Qty & Part & Size & Job \\
\midrule
1 & M10 brass flat washer (McMaster 98691A118) & OD 20, ID 10.5, 2 thick & Outer contact. Measure one, these are the standard sizes \\
1 & M3 hex-head bolt (McMaster 91280A104) & Head 5.5 flats, 2 tall & Center contact. The flat top of the head faces down \\
1 & M3 hex nut & 5.5 flats, 2.4 thick & Clamps one resistor lead to the bolt \\
2 & Disc magnet, KJ Magnetics D52 & \O{} 7.94 $\times$ 3.175 & 5/16'' $\times$ 1/8''. Pulls the piece down onto its pads \\
1 & Resistor, 1/4\,W, 1\% metal film & $\approx$6.3 $\times$ 2.4 body & Tells the board which piece it is (table below) \\
\bottomrule
\end{tabularx}

\subsection{Pockets and the contact plug}

\begin{tabularx}{\textwidth}{l l L}
\toprule
Feature & Size & Notes \\
\midrule
Washer pocket & \O{} 20.3, 2.0 deep & 0.3 of slip on the diameter (\texttt{FIT\_SLIP}) \\
Contact plug (printed) & \O{} 10.2 $\times$ 2.0 & Sits in the washer hole and insulates the bolt from the washer \\
Plug flange & 16 $\times$ 5.5 bar, 0.9 thick & Sits on top of the washer. A bar, not a disc, so the magnets fit beside it \\
Flange pocket & 16.3 $\times$ 5.8, 1.0 deep & \\
Bolt head pocket in plug & Hexagon, 5.8 across flats, 2.0 deep & The head sits flush with the bottom of the plug \\
Nut pocket & 6.65 $\times$ 5.8 box & Tight on the flats, so the nut can't spin while you tighten the bolt \\
Magnet pockets & \O{} 8.24 $\times$ 3.4, at $y = \pm$7.82 & Closed. 0.45 of plastic under them, 0.6 over them \\
Resistor cavity & \O{} 12, 10 tall, starts at $z = 6.4$ & 45\textdegree{} cone roof, no supports \\
Lead slot & 1.6 wide, out to r 9.6 & One resistor lead comes down here and gets soldered to the washer \\
\bottomrule
\end{tabularx}

\subsection{Magnets}

Brass isn't magnetic, so the magnets go in the pieces and grab a steel washer buried under each
square. The steel washer is centered on the square, so it also pulls a piece toward the middle.
There's 3.45\mm{} of plastic and air between the magnets and the steel. \file{checks/fit\_check.py}
checks that every magnet pocket has at least 0.8\mm{} of wall around it.

\begin{itemize}
  \item \textbf{Putting them in:} the pockets are closed, so add a print pause in PrusaSlicer at the
        first layer that covers them, about 5.5\mm{} up from the bed (right-click the \textbf{+} on the
        layer slider, then \emph{Add pause print}). When the printer stops, drop in two magnets and resume.
  \item \textbf{Which way up} doesn't matter, since they're pulling on plain steel.
  \item \textbf{Too strong or too weak?} Test a few pieces on the 2$\times$2 tile first. To weaken the
        pull, raise \texttt{TARGET\_SKIN} in \file{dims.py}. To strengthen it, lower it (0.6 at least).
  \item \textbf{The D62 magnets} (3/8'') are too wide for the pawn. Keep them for pull tests.
\end{itemize}

\subsection{Resistor values}

The divider is 3.3\,V $\rightarrow$ piece resistor $\rightarrow$ sense node $\rightarrow$ 10\,k$\Omega$ to ground.
The 12 levels are about 0.25\,V apart. Even with both resistors off by a full 1\%, neighboring
levels stay at least 0.18\,V apart.

\begin{tabularx}{\textwidth}{c l c c L}
\toprule
Level & Piece & Resistor (E24, 1\%) & Expected voltage & Count for one set \\
\midrule
0  & Empty square & none   & 0.000\,V & \\
1  & White pawn   & 120\,k & 0.254\,V & 8 \\
2  & White knight & 56\,k  & 0.500\,V & 2 \\
3  & White bishop & 33\,k  & 0.767\,V & 2 \\
4  & White rook   & 22\,k  & 1.031\,V & 2 \\
5  & White queen  & 16\,k  & 1.269\,V & 2 (1 + spare for promotion) \\
6  & White king   & 12\,k  & 1.500\,V & 1 \\
7  & Black pawn   & 8.2\,k & 1.813\,V & 8 \\
8  & Black knight & 6.2\,k & 2.037\,V & 2 \\
9  & Black bishop & 4.3\,k & 2.308\,V & 2 \\
10 & Black rook   & 3.0\,k & 2.538\,V & 2 \\
11 & Black queen  & 1.8\,k & 2.797\,V & 2 (1 + spare for promotion) \\
12 & Black king   & 820    & 3.050\,V & 1 \\
\bottomrule
\end{tabularx}

That's 34 pieces in all: 32 for the game plus a spare queen for each color.

% ---------------------------------------------------------------------------
\section{Claw}

One servo turns a pinion. The pinion drives two racks, one on each jaw, so the jaws slide
apart or together along X. The jaws close to a fixed spot under the cap and never squeeze.
When the claw lifts, the cap cone drops into the countersink in the jaw lips and the piece
centers itself. When the claw opens, the piece drops off.

\subsection{Printed parts}

\begin{tabularx}{\textwidth}{c l l L}
\toprule
Qty & File & Overall size & Print orientation \\
\midrule
2 & \file{STL/claw/jaw.stl}     & 37.0 $\times$ 27.1 $\times$ 20.0 & Upside down, flat top of the rack on the bed \\
1 & \file{STL/claw/pinion.stl}  & \O{} 27.5 $\times$ 6.0 & Flat side down, hub recess facing up \\
1 & \file{STL/claw/housing.stl} & 30.0 $\times$ 44.7 $\times$ 14.8 & Upside down, servo-tab face on the bed \\
\bottomrule
\end{tabularx}

The second jaw is the same part turned 180\textdegree.

\subsection{Jaws}

\begin{tabularx}{\textwidth}{l l L}
\toprule
Feature & Size & Notes \\
\midrule
Closed opening (seat) & r 9 & Neck r 5 plus 4 of slop, so closing never touches the neck \\
Seat countersink & 45\textdegree, r 9 to r 12 & Matches the cap cone \\
Lip & 5 thick, 12 wide & $z = -5$ to $0$ \\
Finger & 4 thick, 16 wide & Inner face at $x = 17$, 5 outside the cap \\
Stroke & 10 per jaw & Open jaws clear a cap that is 4 off center \\
Closing height & 5 below seated & The claw closes 5 lower than where the piece ends up \\
Grab tolerance & $\pm$4 & Verified in 8 directions by \file{checks/claw\_check.py} \\
Open footprint & X $\pm$31, Y $\pm$22.4 & Misses a king pushed 4 toward the target on the next square \\
\bottomrule
\end{tabularx}

\subsection{Rack and pinion}

\begin{tabularx}{\textwidth}{l l L}
\toprule
Feature & Size & Notes \\
\midrule
Gear module & 1.25 & 20\textdegree{} pressure angle, involute pinion, straight-sided rack \\
Pinion teeth & 20 & \\
Pinion pitch / tip / root \O & 25.0 / 27.5 / 21.875 & \\
Backlash & 0.15 & Shaved off every tooth so printed gears don't bind \\
Gear thickness & 6 & Rack and pinion both sit at $z = 9$ to $15$ \\
Rack backing & 5 & Solid material behind the tooth roots \\
Rack chamfer & 2.5 at 45\textdegree & Rides on a matching ledge in the housing \\
Servo turn for full stroke & 45.8\textdegree & 10\mm{} / 12.5\mm{} pitch radius. That's 4.58\textdegree{} per mm \\
Pinion bore & \O{} 5.0 through & Reach the horn screw from below \\
Hub recess & \O{} 8.0, 1.5 deep & Fits the servo horn's hub \\
\bottomrule
\end{tabularx}

\subsection{Housing}

\begin{tabularx}{\textwidth}{l l L}
\toprule
Feature & Size & Notes \\
\midrule
Rack guides & $x = -15$ to $15$, walls 3 thick & Run along both racks \\
Sliding clearance & 0.3 & Between racks and guides. Raise it in \file{dims.py} if the jaws bind \\
Horn swing clearance & r 17, 3 tall & Between the pinion and the servo \\
Servo pocket & 12.7 $\times$ 23.3 & 0.25 of slip per side, shifted 5.5 in Y so the shaft is centered \\
Servo screw pilot holes & \O{} 1.8, 6 deep, 27.8 apart & For the M2 screws that come with the servo \\
Heat-set insert holes & 4 $\times$ \O{} 4.0, 5 deep, at ($\pm$11, $\pm$18) & M3 inserts, for bolting to the Z carriage \\
\bottomrule
\end{tabularx}

Heights above the seat plane: jaw bottoms $-5$, rack bottoms 9, rack tops 15, housing bottom 15.3,
servo face 18, housing top (servo tabs sit here) 22.3, servo top $\approx$40.7.

\subsection{Bought parts for the claw}

\begin{tabularx}{\textwidth}{c l L}
\toprule
Qty & Part & Notes \\
\midrule
1 & MG90S micro servo (metal gear) & Comes with horns and M2 screws \\
4 & M3 heat-set inserts & \O{} 4.0 hole, 5 deep. Press in at about 220\,\textdegree C \\
4 & M3 bolts & \todo{length depends on the Z carriage plate} \\
\bottomrule
\end{tabularx}

% ---------------------------------------------------------------------------
\section{Servo (MG90S): measure before printing the housing}

These numbers came from the datasheet. Check your actual servo with calipers. If anything
is off by more than a few tenths, change it in \file{dims.py} and rebuild the housing.

\begin{tabularx}{\textwidth}{l c L}
\toprule
Dimension & Datasheet & Measured \\
\midrule
Case length & 22.8 & \\
Case width & 12.2 & \\
Case height (bottom to output face) & 22.7 & \\
Shaft offset from case center & 5.5 & \\
Output face to underside of tabs & 4.3 & \\
Tab span (end to end) & 32.5 & \\
Tab thickness & 2.5 & \\
Tab width & 5.4 & \\
Mounting hole spacing & 27.8 & \\
\bottomrule
\end{tabularx}

% ---------------------------------------------------------------------------
\section{Gantry and motors}

It's a Cartesian gantry, basically a 3D printer with a claw instead of a nozzle. The frame is
the aluminum extrusion we already have. Four 2525 posts stand at the corners, just outside the
board. A printed corner block sits on each post, and the two 600\mm{} 3030 runways rest on the
corner blocks, one along each side of the board. An MGN12 linear rail is bolted on top of each
runway, and the 450\mm{} 3030 bridge sits on the two rail carriages, so it rolls along the board
in X. A third rail on the side of the bridge carries the claw in Y. The models are in
\file{cad/src/gantry/}, and \file{checks/gantry\_check.py} checks the travel and clearances.
The step-by-step build is in \file{gantry-guide.pdf}.

\subsection{Travel the gantry needs}

\begin{tabularx}{\textwidth}{l l L}
\toprule
Axis & Range & Why \\
\midrule
X & $\pm$275 & Squares reach $\pm$175. Capture columns sit at $\pm$225 and $\pm$275, one side per color \\
Y & $\pm$175 & Squares and capture slots (2 columns of 8 per side) \\
Z (seat height) & 38 to 169 & Lowest: closing on a pawn ($45 - 2 - 5$). Highest: carrying a king over another king with 5 to spare \\
Z travel & 131 & \\
\bottomrule
\end{tabularx}

\subsection{Frame (board coordinates)}

\begin{tabularx}{\textwidth}{l l L}
\toprule
Part & Size and place & Notes \\
\midrule
Runways & 2 $\times$ 3030, 600 long & Centers at $y = \pm$210, bottom at $z = 230$, from $x = -360$ to 240. Used as is \\
Bridge & 3030, 450 long & Bottom at $z = 279$. Its ends sit right over the runways. Used as is \\
Posts & 4 $\times$ 2525, \textbf{cut to 231} & At $x = -344.3$ and 224.3, $y = \pm$242. All four the same length, square ends \\
Corner blocks & Printed, 4 & The runway rests on a 6\mm{} saddle, the post goes 30\mm{} up into a sleeve \\
Post feet & Printed, 4 & Screwed to the plywood base. The posts stand 4.5\mm{} outside the board \\
Bridge brackets & Printed, 2 & Bolt to the X carriages, hold the bridge, and clamp the X belts \\
Plywood base & About 720 $\times$ 605 & The board sits on it too. Board top ($z = 0$) is 10 above it \\
Whole frame & 676 $\times$ 591, 319 tall & Not centered on the board in X: the motors stick out past the $-X$ end \\
\bottomrule
\end{tabularx}

The claw hangs 60\mm{} in front of the bridge (\texttt{CLAW\_OFFSET\_X}), so the runways are shifted
60 toward $-X$ to keep the claw's reach centered. \todo{60 is a guess until the Z carriage is designed.}

\subsection{Rails and clearances}

\begin{tabularx}{\textwidth}{l l L}
\toprule
Check & Value & Notes \\
\midrule
X rails & 2 $\times$ MGN12, 600, MGN12\textbf{C} carriage & The short C carriage gives 565 of travel. We need 550, so 7.6 to spare at each end \\
Y rail & MGN12, 450, MGN12H carriage & 404 of travel, 27 to spare at each end \\
Claw under the runways & Claw top at 216, runways at 230 & With a king hanging from the claw at its highest \\
Claw beside the posts & 29 to spare & At the corner squares, jaws open \\
Frame vs.\ board & 1.3 gap & Sleeves and feet stay outside the board, so it lifts straight out \\
\bottomrule
\end{tabularx}

\subsection{Motors and drive}

\begin{tabularx}{\textwidth}{c l L}
\toprule
Qty & Part & Notes \\
\midrule
4 & NEMA 17, 42\,Ncm, 1.5\,A (17HE15-1504S) & Two for X (one per runway), one for Y, one for Z \\
1 & BTT SKR Mini E3 V3 + Marlin & Four TMC2209 drivers. Both X motors share the X driver through a splitter cable \\
2 & GT2 belt, 6 wide, about 1400 each & One loop per runway. \todo{Y belt length depends on the Y motor mount} \\
3 & GT2 20-tooth pulley, 5\mm{} bore & 40\mm{} of belt per turn, so 80 steps/mm at 16 microsteps \\
3 & GT2 toothless idler, 5\mm{} bore & Two for X, on an M5 bolt in the idler block's slot. One for Y \\
1 & T8 lead screw, 8\mm{} lead & 400 steps/mm at 16 microsteps. Holds position with the power off \\
3 & Limit switches & Home X to $-X$, Y to $-Y$, Z to the top. \todo{switch mounts aren't designed yet} \\
1 & 24\,V power supply, 150\,W & \\
\bottomrule
\end{tabularx}

The X belt runs in a loop just outside each runway: motor pulley at $x = -385$, idler at $x = 255$,
both at $z = 262$. The top strand is clamped to the bridge bracket. The two X motors face opposite
ways, so \textbf{one of them has to be wired backwards} or they'll twist the bridge. The guide
covers this.

\subsection{Frame hardware}

\begin{tabularx}{\textwidth}{c l L}
\toprule
Qty & Part & Where \\
\midrule
8 & M6 $\times$ 12 + M6 tee nut (8\mm{} slot) & Corner block saddles up into the runways \\
4 & M6 $\times$ 10 + M6 tee nut & Bridge bracket cheeks into the bridge \\
12 & M5 $\times$ 10 + M5 tee nut (6\mm{} slot) & Corner block sleeves and post feet into the posts \\
8 & M3 $\times$ 6 & Bridge brackets down into the X carriages \\
8 & M3 $\times$ 10 & X motors to the corner blocks \\
about 35 & M3 $\times$ 10 + M3 tee nut (8\mm{} slot) & MGN12 rails onto the extrusion, every other hole \\
2 & M5 $\times$ 30 bolt, nut, washers & X idlers \\
16 & \#8 $\times$ 3/4'' wood screws & Post feet into the base \\
\bottomrule
\end{tabularx}

Check the slot widths on your extrusion first. The tee nut sizes above assume an 8\mm{} slot on
the 3030 and a 6\mm{} slot on the 2525.

\todo{the Y motor mount, the Y carriage and Z axis, and the limit switch mounts. The claw's
placeholder in the assembly shows where they go.}

% ---------------------------------------------------------------------------
\section{Board sensing electronics}

\begin{tabularx}{\textwidth}{c l L}
\toprule
Qty & Part & Notes \\
\midrule
64 & 10\,k$\Omega$ 1\% reference resistor & One per square, so the mux's own resistance stays out of the divider \\
4 & CD74HC4067 16:1 analog mux & 64 squares into one ADC channel. A full scan takes about 1\,ms \\
1 & ADC & \todo{which part. 10-bit is enough. 12-bit gives about 220 counts between levels} \\
1 roll & Copper tape, at least 26 wide & For the pads on each square \\
64 & M10 zinc-plated steel washer & Glued under each square for the piece magnets to grab \\
& Hookup wire & Through the two wire holes on each square \\
\bottomrule
\end{tabularx}

To decode a reading, snap it to the nearest level. If it's more than about 0.09\,V from every
level, flag the square as unreadable and don't guess.

% ---------------------------------------------------------------------------
\section{Printing (Prusa MK4, PLA)}

\begin{tabularx}{\textwidth}{l L}
\toprule
Setting & Value \\
\midrule
Printer profile & Original Prusa MK4 Input Shaper, 0.4 nozzle \\
Filament & Generic PLA \\
Print profile & 0.15\,mm QUALITY \\
Perimeters & 4 \\
Infill & 25\% gyroid \\
Supports & None (every part was designed to print without them) \\
Brim & Off unless a corner lifts \\
Elephant foot compensation & Default (0.2). The sliding faces and the servo pocket are on the bed \\
\bottomrule
\end{tabularx}

The whole claw (2 jaws, pinion, housing) fits on one plate. It's about 16\,cm$^3$ of plastic
and roughly 1.5--2 hours of printing.

\textbf{Pieces:} 100\% infill so they have some weight, standing up, with a print pause at about
5.5\mm{} to drop in the magnets (see Magnets above). The contact plugs print flange side down.

\textbf{Gantry parts:} 4 perimeters and 40\% infill, since they carry the frame. Each STL in
\file{STL/gantry/} is already laid in its print orientation: the corner blocks on their outside face,
the bridge brackets on their cheek, the post feet on their flange. The \file{plus\_y} and
\file{minus\_y} versions are mirror images, so print one of each.

\end{document}
